% Generated draft: requires mathematical and source review.
\documentclass[11pt]{article}
\usepackage[margin=25mm]{geometry}
\usepackage{fontspec}
\setmainfont{Latin Modern Roman}
\usepackage{amsmath,amssymb,mathtools,unicode-math}
\setmathfont{Latin Modern Math}
\usepackage{enumitem,needspace,xurl,hyperref}
\setlength{\parindent}{0pt}
\setlength{\parskip}{5pt}
\setlength{\emergencystretch}{3em}
\newcommand{\N}{\mathbb{N}}
\newcommand{\Z}{\mathbb{Z}}
\newcommand{\Q}{\mathbb{Q}}
\newcommand{\R}{\mathbb{R}}
\newcommand{\C}{\mathbb{C}}
\newcommand{\F}{\mathbb{F}}
\newcommand{\A}{\mathbb{A}}
\newcommand{\PP}{\mathbb{P}}
\newcommand{\E}{\mathbb{E}}
\newcommand{\eps}{\varepsilon}
\newcommand{\dd}{\,\mathrm{d}}
\begin{document}
\section*{UnsolvedMath ID: 30000594}
\textbf{Source problem code:} catalogue entry
\subsection{Definitions and mathematical statement}
\textbf{Index-comparison problem.} For a closed smooth spin manifold \(M\), what is the relation between its Rosenberg indices \(\alpha_{\mathrm{red}}(M)\), \(\alpha_{\mathrm{max}}(M)\), and the codimension-two positive-scalar-curvature obstructions of Gromov and Lawson? The source asks for a relation, without specifying one equivalence or implication as the conjectured target.\par In the usual connected real-spin convention, let \(n=\dim M\), \(\pi=\pi_1(M)\), and let \(C^*_{\mathrm{max},\R}\pi\) and \(C^*_{\mathrm r,\R}\pi\) denote the maximal and reduced real group \(C^*\)-algebras. The Rosenberg classes are the higher indices of the spin Dirac operator twisted by the corresponding flat Mishchenko bundles \(\widetilde M\times_\pi C^*\pi\):\[\alpha_{\mathrm{max}}(M)\in KO_n(C^*_{\mathrm{max},\R}\pi),\qquad \alpha_{\mathrm{red}}(M)\in KO_n(C^*_{\mathrm r,\R}\pi).\] The canonical quotient of group \(C^*\)-algebras sends the maximal class to the reduced class. In particular reduced nonvanishing entails maximal nonvanishing, while the converse is not asserted merely from this map. A positive-scalar-curvature metric forces both classes to vanish by the spin Dirac obstruction. The source does not specify real versus complex index conventions, so the displayed real formulation is explicit rather than silently interchangeable with complex \(K\)-theory.\par The codimension-two obstructions concern suitable embedded submanifolds \(N\subset M\) of codimension two, with additional normal-bundle, fundamental-group, topology, and nonvanishing or geometric hypotheses. The source does not quote the precise Gromov--Lawson criterion or fix those data. The problem includes determining whether and under which hypotheses such an obstruction is detected by a nonzero ambient Rosenberg class, and what comparison mechanism relates the codimension-two data to the ambient indices. Any degree-lowering transfer map would require its own construction and assumptions; it is not supplied automatically by an inclusion of fundamental groups. No general equivalence between nonexistence of positive scalar curvature and nonvanishing of either index is asserted.\par The supplied triage checked 2026-08-21 labels the item partially solved. It reports substantial development of codimension-two index obstructions and their relation to Rosenberg classes, but says the general relation is not completely settled. The cited theorems and their exact hypotheses are not reproduced. This dated assessment requires independent verification.
\subsection{Short English statement}
Determine which codimension-two curvature obstructions are detected by the ambient maximal or reduced higher Dirac index, and under what comparison hypotheses. The precise submanifold criterion is missing. The standard maximal-to-reduced index map alone does not supply every requested relation.
\subsection{Sources}
\begin{itemize}\item Local UnsolvedMath registry, permanent ID 30000594, code catalogue entry; \emph{Rosenberg Indices and Codimension-Two Curvature Obstructions}, Oberwolfach Reports Open Problems; original attribution to Gromov and Lawson for the codimension-two obstructions.\item \emph{Topologie (2007)}, Oberwolfach Reports, DOI 10.4171/owr/2006/43, pp. 2579--2630. \url{https://doi.org/10.4171/owr/2006/43}. Proposed year and year tag: 2006; no proposer named.\item Supplied literature label \emph{original OWR report}: \url{https://doi.org/10.4171/owr/2006/43}.\item Supplied literature label \emph{Hanke--Pape--Schick, codimension-two obstructions}: \url{https://doi.org/10.5802/aif.3000}.\item Supplied literature label \emph{Kubota--Schick, C-star index}: \url{https://doi.org/10.2140/gt.2021.25.949}.\item Literature triage checked 2026-08-21, status \texttt{partially\_solved}; record updated 2026-09-16. The \texttt{corrected\_verified} field records notation and punctuation changes only. Verify dated claims independently.\end{itemize}
\end{document}
